\chapter{State of the Art}
\section{Hybrid semantic query languages}
SUQL extends SQL with \textsc{Answer} and \textsc{Summary} functions over free-text
columns and translates conversational input using in-context learning
\cite{liu2024suql}. Its representation is expressive, but a physical executor can
still invoke an NLP model once per tuple. BlendSQL similarly embeds external
reasoning in relational algebra \cite{glenn2024blendsql}. This project operates at
the physical-plan level: it does not propose a new language.

\section{Semantic joins and batch prompting}
Trummer's semantic block join is inspired by block nested-loop join and includes
multiple rows from both inputs in one prompt \cite{trummer2025semanticjoins}. Batch
prompting amortizes fixed instructions and request overhead and can reduce time and
token cost, but large batches risk omitted identifiers and interference
\cite{cheng2023batch}. Our join key is not semantic: movie and review identifiers
must be exactly equal. Only the residual review predicate belongs in the LLM.

\section{Model cascades}
Model cascades let a cheap model answer easy cases and escalate difficult cases.
FrugalGPT demonstrates that learned cascades can reduce inference cost while
maintaining quality \cite{chen2023frugalgpt}. In this project, a large model labels
an online calibration sample because no per-predicate labelled training set is
available at execution time. This estimates agreement with the large model, not
ground-truth correctness; the distinction is central to the limitations.

\section{Positioning}
The proposed plan composes three orthogonal ideas. Relational selection reduces the
candidate cardinality, batching reduces requests per candidate, and cascading
reduces the fraction handled by the large model. The comparison includes complete
end-to-end systems, so a future factorial ablation is needed to attribute gains
causally to individual components.
